% Page design of the PDF of the ESB documentation (XeLaTeX): colors, fonts, headings, header and footer, lists, tables, and figures.
% docs/conf.py puts this file into the key preamble of latex_elements, after \documentid, \versionid, and \esbtitle, and after
% the packages and the options of \sphinxsetup of Sphinx, so that what is set here replaces them.

\definecolor{esbdark}{HTML}{1D3546}
\definecolor{esborange}{HTML}{F0690A}

\usepackage{eso-pic}     % the background of the title page
\usepackage{sectsty}
\usepackage{titlesec}
\usepackage{fancyhdr}
\usepackage{enumitem}
\usepackage{etoolbox}
\usepackage{array}
\usepackage{caption}
\usepackage[strings]{underscore}
% a long identifier may break after an underscore
\renewcommand{\_}{\textunderscore\allowbreak}

% -- Fonts ---------------------------------------------------------------------------------------------------------------
% Inter for the text and Inconsolata for the code, which the packages fonts-inter and fonts-inconsolata install. Where
% one is missing, a font that is common to systems takes its place.
\usepackage{fontspec}
\IfFontExistsTF{Inter}{\newcommand{\esbtextfont}{Inter}}{%
  \IfFontExistsTF{Arial}{\newcommand{\esbtextfont}{Arial}}{%
    \IfFontExistsTF{Liberation Sans}{\newcommand{\esbtextfont}{Liberation Sans}}{\newcommand{\esbtextfont}{DejaVu Sans}}}}
\IfFontExistsTF{Inconsolata}{\newcommand{\esbcodefont}{Inconsolata}}{%
  \IfFontExistsTF{Courier New}{\newcommand{\esbcodefont}{Courier New}}{%
    \IfFontExistsTF{Liberation Mono}{\newcommand{\esbcodefont}{Liberation Mono}}{\newcommand{\esbcodefont}{DejaVu Sans Mono}}}}
% fontspec takes Inter ExtraBold for the bold face of Inter, Bold is the one that is meant. The contextual alternates of Inter
% (calt) are off: they put a bracket or a hyphen next to a capital into a glyph without a Unicode value, so that the text of
% the PDF could not be copied or searched.
\IfFontExistsTF{Inter Bold}{%
  \setmainfont{\esbtextfont}[Ligatures=TeX, RawFeature={-calt}, BoldFont={Inter Bold}, BoldItalicFont={Inter Bold Italic}]
  \setsansfont{\esbtextfont}[Ligatures=TeX, RawFeature={-calt}, BoldFont={Inter Bold}, BoldItalicFont={Inter Bold Italic}]
}{%
  \setmainfont{\esbtextfont}[Ligatures=TeX, RawFeature={-calt}]
  \setsansfont{\esbtextfont}[Ligatures=TeX, RawFeature={-calt}]
}
\setmonofont{\esbcodefont}[Scale=MatchLowercase, AutoFakeBold=2]

% -- Text ----------------------------------------------------------------------------------------------------------------
% A word moves to the next line instead of breaking.
\hyphenpenalty=10000
\exhyphenpenalty=10000

% no blank page before a chapter
\makeatletter
\@openrightfalse
\makeatother

\addto\captionsenglish{\renewcommand{\contentsname}{Table of contents}}

% inline code a little smaller than the text, and the size of the code blocks
\AtBeginDocument{\renewcommand{\sphinxcode}[1]{\texttt{\small #1}}}
\fvset{fontsize=\small}

% links are blue, except the ones of the table of contents
\hypersetup{colorlinks=true, linkcolor=blue, citecolor=blue, urlcolor=blue, filecolor=blue, linktoc=all}
\let\esbtableofcontents\tableofcontents
\renewcommand{\tableofcontents}{{\hypersetup{linkcolor=black}\esbtableofcontents}}

% -- Headings ------------------------------------------------------------------------------------------------------------
\allsectionsfont{\color{esbdark}}

\titleformat{\chapter}[display]
  {\normalfont\huge\bfseries\color{esbdark}\filcenter}
  {\chaptertitlename\ \thechapter}
  {2pt}
  {\Huge}
  [\vspace{-8pt}{\color{esborange}\rule{\textwidth}{0.8pt}}]
\titlespacing*{\chapter}{0pt}{-20pt}{12pt}

\titleformat{\section}{\normalfont\Large\bfseries\color{esbdark}}{\thesection}{1em}{}
\titlespacing*{\section}{0pt}{3.5ex plus 1ex minus .2ex}{2.3ex plus .2ex}
\titleformat{\subsection}{\normalfont\large\bfseries\color{esbdark}}{\thesubsection}{1em}{}
\titlespacing*{\subsection}{0pt}{3.25ex plus 1ex minus .2ex}{1.5ex plus .2ex}
\titleformat{\subsubsection}{\normalfont\normalsize\bfseries\color{esbdark}}{\thesubsubsection}{1em}{}
\titlespacing*{\subsubsection}{0pt}{3.25ex plus 1ex minus .2ex}{1.5ex plus .2ex}

% -- Header and footer ---------------------------------------------------------------------------------------------------
% The title in the header, the page number below an orange rule in the footer; the first page of a chapter has them too.
\setlength{\headheight}{14pt}
\newcommand{\esbpagestyle}{%
  \fancyhf{}%
  \fancyhead[C]{\small\bfseries\esbtitle}%
  \fancyfoot[C]{\small\thepage}%
  \renewcommand{\headrulewidth}{0pt}%
  \renewcommand{\footrulewidth}{0.8pt}%
  \renewcommand{\footrule}{{\color{esborange}\hrule width\headwidth height\footrulewidth}}%
}
\fancypagestyle{normal}{\esbpagestyle}
\fancypagestyle{plain}{\esbpagestyle}
\pagestyle{normal}

% -- Lists ---------------------------------------------------------------------------------------------------------------
\newlength{\esbindent}
\setlength{\esbindent}{2em}

\setlist[itemize,1]{leftmargin=\dimexpr\esbindent+1.25em\relax, label=\textbullet, labelsep=0.5em, labelwidth=1em, align=parleft,
  itemsep=2pt, parsep=0pt, topsep=4pt, partopsep=2pt, listparindent=0pt, itemindent=0pt}
\setlist[enumerate,1]{leftmargin=\dimexpr\esbindent+1.25em\relax, labelsep=0.5em, labelwidth=1em, align=parleft,
  itemsep=2pt, parsep=0pt, topsep=4pt, partopsep=2pt, listparindent=0pt, itemindent=0pt}
\setlist[itemize,2,3]{leftmargin=\dimexpr\esbindent+2.75em\relax, labelsep=0.25em, labelwidth=1em, align=parleft,
  itemsep=1pt, parsep=0pt, topsep=2pt, listparindent=0pt, itemindent=0pt}
\setlist[enumerate,2,3]{leftmargin=\dimexpr\esbindent+2.75em\relax, labelsep=0.25em, labelwidth=1em, align=parleft,
  itemsep=1pt, parsep=0pt, topsep=2pt, listparindent=0pt, itemindent=0pt}

% -- Tables and figures --------------------------------------------------------------------------------------------------
% Tables are small, start at the left margin, and have bold column titles.
\tymin=15pt
\tymax=\maxdimen
\setlength{\arrayrulewidth}{0.8pt}
\AtBeginDocument{%
  \let\esbattablestart\sphinxattablestart
  \renewcommand{\sphinxattablestart}{%
    \par\medskip
    \let\centering\raggedright
    \raggedright
    \esbattablestart
    \footnotesize
    \setlength{\tabcolsep}{2pt}%
    \setlength{\emergencystretch}{3em}%
  }%
  \def\sphinxstyletheadfamily{\bfseries}%
}
\AtBeginEnvironment{longtable}{%
  \setlength{\LTleft}{0pt}%
  \setlength{\LTright}{0pt}%
  \footnotesize
  \setlength{\tabcolsep}{2pt}%
  \setlength{\emergencystretch}{3em}%
}

\captionsetup[figure]{position=bottom, labelfont=it, textfont=it, labelsep=period, name={Fig.~}, justification=centering,
  singlelinecheck=false}

% A PNG figure can carry a resolution in the file, which LaTeX honors and browsers ignore; every figure is as wide as the text.
\setkeys{Gin}{width=\linewidth}
